% ============================================================================
%  3.2 DEFINICIÓN DE LA FUNCIONALIDAD
% ============================================================================

\section{Definici\'on de la funcionalidad}

\subsection{Nombre de la funcionalidad}

Inscripci\'on a Eventos Deportivos.

\subsection{Objetivo}

Permitir a los usuarios autenticados inscribirse a eventos deportivos
asociados a los escenarios, cancelar dichas inscripciones, y consultar
su estado desde la interfaz de la aplicaci\'on.

\subsection{Usuarios involucrados}

\begin{table}[H]
\centering
\begin{tabular}{@{}lll@{}}
\toprule
\textbf{Rol} & \textbf{Acciones} & \textbf{Restricciones} \\
\midrule
Asistente & Inscribirse, cancelar, ver sus inscripciones & Solo sus propias \\
Gestor & Ver inscripciones de sus eventos & Solo sus escenarios \\
Admin & Ver todas las inscripciones, eliminar cualquiera & Sin restricciones \\
\bottomrule
\end{tabular}
\caption{Roles y permisos del m\'odulo}
\end{table}

\subsection{Flujo de funcionamiento}

El flujo principal de la funcionalidad se describe a continuaci\'on:

\begin{figure}[H]
\centering
\begin{tikzpicture}[
  node distance=0.8cm,
  caja/.style={rectangle, draw=black!60, rounded corners,
    minimum width=2.8cm, minimum height=0.7cm, align=center,
    font=\footnotesize, fill=blue!5},
  flecha/.style={-Stealth, thick}
]
  \node[caja] (inicio) {Usuario};
  \node[caja, right=of inicio] (escenario) {Detalle del\\escenario};
  \node[caja, right=of escenario] (evento) {Selecciona\\evento};
  \node[caja, below=of evento] (validar) {Validar\\datos};
  \node[caja, left=of validar] (guardar) {Guardar en\\Supabase};
  \node[caja, left=of guardando] (exito) {Inscripci\'on\\confirmada};
  \node[caja, below=of validar] (error) {Mostrar\\error};

  \draw[flecha] (inicio) -- (escenario);
  \draw[flecha] (escenario) -- (evento);
  \draw[flecha] (evento) -- (validar);
  \draw[flecha] (validar) -- node[above, font=\scriptsize] {V\'alido} (guardar);
  \draw[flecha] (validar) -- node[right, font=\scriptsize] {Inv\'alido} (error);
  \draw[flecha] (guardar) -- (exito);
\end{tikzpicture}
\caption{Flujo de inscripci\'on a eventos}
\end{figure}

\subsection{Operaciones CRUD}

El m\'odulo implementa las siguientes operaciones sobre la tabla
\texttt{event\_registrations}:

\begin{table}[H]
\centering
\begin{tabular}{@{}llp{6cm}@{}}
\toprule
\textbf{Operaci\'on} & \textbf{M\'etodo HTTP} & \textbf{Descripci\'on} \\
\midrule
CREATE & INSERT & El usuario se inscribe a un evento \\
READ & SELECT & Consultar inscripciones propias o de un evento \\
UPDATE & UPDATE & Cambiar estado de inscripci\'on (no implementado
  actualmente, se usa DELETE) \\
DELETE & DELETE & Cancelar inscripci\'on \\
\bottomrule
\end{tabular}
\caption{Operaciones CRUD del m\'odulo}
\end{table}

\subsection{Persistencia}

Toda la informaci\'on de inscripciones se almacena en Supabase
(PostgreSQL), ya que requiere persistencia centralizada:
\begin{itemize}
  \item \textbf{Inspecci\'on:} Disponible para todos los dispositivos
    del usuario (multi-dispositivo).
  \item \textbf{Gesti\'on:} Los gestores y admins pueden consultar
    inscripciones desde cualquier punto.
  \item \textbf{Seguridad:} Las pol\'iticas RLS garantizan que cada
    usuario solo vea sus propias inscripciones (salvo admin/gestor).
\end{itemize}
